\newcommand{\langue}{english}
\newcommand{\nomchapitre}{Chapter}
\newcommand{\titrecourt}{REFERENCE MANUAL}
\input{version.tex}
% ================================================================
%  Préambule du manuel de Coupole (FR et EN), d'après le format
%  maison : Latin Modern, titres sans-serif colorés, encadrés
%  tcolorbox, fontawesome5.  \langue{french} ou \langue{english}
%  est défini par le fichier principal avant \input.
% ================================================================
\documentclass[11pt,a4paper]{article}

\usepackage[T1]{fontenc}
\usepackage[utf8]{inputenc}
\usepackage[\langue]{babel}
\usepackage{lmodern}
\usepackage{textcomp}
\usepackage{amsmath,amssymb}
\usepackage{mathtools}
\usepackage{siunitx}
\sisetup{group-minimum-digits=4,separate-uncertainty=true,per-mode=symbol}
\usepackage{cancel}
\usepackage{microtype}
\DeclareUnicodeCharacter{2248}{\ensuremath{\approx}}
\DeclareUnicodeCharacter{2212}{\ensuremath{-}}
\DeclareUnicodeCharacter{00D7}{\ensuremath{\times}}
\DeclareUnicodeCharacter{2192}{\ensuremath{\rightarrow}}
\DeclareUnicodeCharacter{00A0}{~}
\DeclareUnicodeCharacter{202F}{\,}
\DeclareUnicodeCharacter{2033}{\ensuremath{''}}
\DeclareUnicodeCharacter{2032}{\ensuremath{'}}
\DeclareUnicodeCharacter{2264}{\ensuremath{\le}}
\DeclareUnicodeCharacter{2265}{\ensuremath{\ge}}
\DeclareUnicodeCharacter{2194}{\ensuremath{\leftrightarrow}}
\DeclareUnicodeCharacter{00B2}{\textsuperscript{2}}
\DeclareUnicodeCharacter{2026}{\ldots}
\DeclareUnicodeCharacter{00B0}{\textdegree}

\usepackage[a4paper,margin=2cm,headheight=24pt]{geometry}
\usepackage{graphicx}
\graphicspath{{img/}}
\usepackage[table]{xcolor}
\usepackage[most]{tcolorbox}
\usepackage{booktabs}
\usepackage{tabularx}
\usepackage{array}
\usepackage{multicol}
\usepackage{enumitem}
\usepackage[nobottomtitles*]{titlesec}
\usepackage{needspace}
\usepackage{fancyhdr}
\usepackage{fontawesome5}
\usepackage{tikz}
\usetikzlibrary{arrows.meta,calc,decorations.pathreplacing,patterns,angles,quotes}
\usepackage{pgfplots}
\pgfplotsset{compat=1.18}
\usepgfplotslibrary{fillbetween}
\usepackage{pdfpages}
\usepackage[colorlinks=true,linkcolor=sky,urlcolor=sky,bookmarksopen=true,bookmarksopenlevel=1,
  pdftitle={Coupole},
  pdfauthor={ARP273-ROSE}]{hyperref}
\usepackage{bookmark}
% Les notes posées dans des tableaux (annexes) par \footnotemark/\footnotetext cassent les ancres
% de hyperref : les appels de note ne sont donc pas des liens.
\hypersetup{hyperfootnotes=false}
\usepackage{xurl}
\usepackage{fvextra}
\fvset{breaklines=true,fontsize=\small}

% ---------------------------------------------------------------- Langue des titres d'encadrés
\newcommand{\tr}[2]{\ifdefstring{\langue}{english}{#2}{#1}}

% ---------------------------------------------------------------- Couleurs
\definecolor{sky}{HTML}{1C6DB4}
\definecolor{skyl}{HTML}{E7F0F9}
\definecolor{alert}{HTML}{B5382B}
\definecolor{alertl}{HTML}{FBEAE7}
\definecolor{night}{HTML}{1B2A4A}
\definecolor{ink}{HTML}{1F2430}
\definecolor{summit}{HTML}{6B7A8F}
\definecolor{gold}{HTML}{B8860B}
\definecolor{goldl}{HTML}{FBF3DC}
\definecolor{pine}{HTML}{2C7A55}
\definecolor{pinel}{HTML}{E2F1E9}
\definecolor{prune}{HTML}{5B3F8C}
\definecolor{prunel}{HTML}{EEE8F6}
\AtBeginDocument{\color{ink}}

% ---------------------------------------------------------------- Titres
% \section = chapitre du cours (« Chapitre 1 : ... »), \subsection = I, II, III..., \subsubsection = A, B, C...
\titleformat{\section}[display]
  {\sffamily\color{night}}
  {\large\bfseries\color{alert}\MakeUppercase{\sectionname}~\thesection}{2pt}{\LARGE\bfseries}
  [{\vspace{2pt}\color{night!35}\titlerule[1.3pt]}]
\titlespacing*{\section}{0pt}{6pt}{10pt}
\newcommand{\sectionname}{\nomchapitre}
\renewcommand{\thesubsection}{\arabic{section}.\arabic{subsection}}
\renewcommand{\thesubsubsection}{}
\titleformat{\subsection}
  {\large\bfseries\sffamily\color{sky}}{\thesubsection}{0.6em}{}
\titlespacing*{\subsection}{0pt}{12pt}{5pt}
\titleformat{\subsubsection}
  {\normalsize\bfseries\sffamily\color{alert}}{}{0em}{}
\titlespacing*{\subsubsection}{0pt}{9pt}{3pt}
\setcounter{secnumdepth}{3}
\setcounter{tocdepth}{2}

\usepackage{etoolbox}
\pretocmd{\section}{\clearpage}{}{}
\pretocmd{\subsection}{\needspace{8\baselineskip}}{}{}
\pretocmd{\subsubsection}{\needspace{5\baselineskip}}{}{}
\renewcommand{\bottomtitlespace}{0.15\textheight}
\widowpenalty=10000
\clubpenalty=10000
\raggedbottom

% ---------------------------------------------------------------- En-têtes
\pagestyle{fancy}
\fancyhf{}
\renewcommand{\headrulewidth}{0.4pt}
\renewcommand{\headrule}{\hbox to\headwidth{\color{night!40}\leaders\hrule height \headrulewidth\hfill}}
\fancyhead[L]{\sffamily\small\color{night}\textbf{COUPOLE · \titrecourt}}
\fancyhead[R]{\sffamily\small\color{summit}\partiecourante}
\fancyfoot[L]{\sffamily\scriptsize\color{summit}Coupole \versioncoupole{} · GPL-3.0 · ARP273-ROSE}
\fancyfoot[R]{\sffamily\small\color{summit}\thepage}
\newcommand{\partiecourante}{}
\newcommand{\partie}[1]{\clearpage\renewcommand{\partiecourante}{#1}}


% ---------------------------------------------------------------- Parties
\titleformat{\part}[display]{\centering\sffamily\color{night}}
  {\Large\bfseries\color{alert}PARTIE \thepart}{8pt}{\Huge\bfseries}
\pretocmd{\part}{\clearpage}{}{}
\setcounter{tocdepth}{2}

% ---------------------------------------------------------------- Encadrés du site (une icône par type de section)
\tcbset{maison/.style={enhanced,breakable,boxrule=0.9pt,arc=2.5pt,left=2.5mm,right=2.5mm,top=1.5mm,bottom=1.5mm,
  fonttitle=\bfseries\sffamily\small,coltitle=white,
  before upper={\widowpenalties 4 10000 10000 10000 0 \clubpenalties 4 10000 10000 10000 0}}}
% \typesite{env}{icône}{libellé}{fond}{cadre}
\newcommand{\typesite}[5]{\newtcolorbox{#1}[1][]{maison,colback=#4,colframe=#5,title={#2~~#3\ifx\relax##1\relax\else~·~##1\fi}}%
  \BeforeBeginEnvironment{#1}{\par\needspace{6\baselineskip}}}
\typesite{cdefinition}{\faBook}{\tr{Définition}{Definition}}{skyl}{sky}
\typesite{crappel}{\faUndo}{\tr{Rappel}{Reminder}}{goldl}{gold!85!black}
\typesite{cremarque}{\faInfoCircle}{\tr{Remarque}{Note}}{skyl}{sky!70!black}
\typesite{cattention}{\faExclamationTriangle}{\tr{Attention}{Warning}}{alertl}{alert}
\typesite{cexemple}{\faLightbulb}{\tr{Exemple}{Example}}{pinel}{pine}
\typesite{cexercice}{\faPencil*}{\tr{Exercice}{Exercise}}{white}{pine}
\typesite{cdemonstration}{\faCalculator}{\tr{Démonstration}{Proof}}{prunel}{prune}
\typesite{cconseil}{\faThumbsUp}{\tr{Conseil}{Tip}}{pinel}{pine!80!black}
\typesite{ccomplement}{\faPlusCircle}{\tr{Complément}{Supplement}}{prunel}{prune!80!black}
\typesite{censavoirplus}{\faSearchPlus}{\tr{En savoir plus}{Further reading}}{prunel}{prune!70!black}
\typesite{cactivite}{\faFlask}{\tr{Activité}{Activity}}{goldl}{gold}
\typesite{cqcm}{\faCheckSquare}{\tr{QCM}{Quiz}}{white}{sky}
\typesite{cobjectifs}{\faBullseye}{\tr{Objectifs}{Goals}}{skyl}{night}
\typesite{cprerequis}{\faListUl}{\tr{Prérequis}{Prerequisites}}{skyl}{night}
\typesite{ccadre}{\faStar}{\tr{À retenir}{Key point}}{goldl}{gold}
\newtcolorbox{ccode}{maison,colback=night!4,colframe=night!50,fontupper=\small\ttfamily,title={\faTerminal~~\tr{Commande}{Command}}}

\newcommand{\intertitre}[2][]{\par\needspace{4\baselineskip}\medskip\noindent{\sffamily\bfseries\color{night}#2}\par\smallskip}
\newcommand{\question}[1]{\par\smallskip\noindent{\sffamily\bfseries\color{pine}#1}~}
\newcommand{\auteurs}[1]{\par\noindent{\sffamily\small\color{summit}\faUserEdit~#1}\par\medskip}
\newcommand{\appli}{\textcolor{summit}{\faMousePointer~\textit{(application interactive en ligne)}}}
% \figsite{titre}{image}{légende}{crédit}
\newcommand{\figsite}[4]{\par\needspace{8\baselineskip}\begin{center}\begin{minipage}{\linewidth}\centering
  \if\relax\detokenize{#1}\relax\else{\sffamily\small\bfseries\color{sky}#1}\par\smallskip\fi
  \if\relax\detokenize{#2}\relax{\sffamily\small\color{summit}\faFilm~\textit{animation ou vidéo en ligne, non reproduite}}\par\else#2\par\fi
  \if\relax\detokenize{#3}\relax\else\smallskip{\small #3}\par\fi
  \if\relax\detokenize{#4}\relax\else{\scriptsize\color{summit}Crédit : #4}\par\fi
  \end{minipage}\end{center}}
\DeclareUnicodeCharacter{2009}{\,}
\DeclareUnicodeCharacter{2013}{--}
\DeclareUnicodeCharacter{2014}{---}
\DeclareUnicodeCharacter{2019}{'}
\DeclareUnicodeCharacter{2018}{`}
\DeclareUnicodeCharacter{201C}{``}
\DeclareUnicodeCharacter{201D}{''}
% \ressource{titre}{niveau · format · volume}{url}{description} — guide du portail
\newcommand{\ressource}[4]{\par\needspace{5\baselineskip}\noindent{\sffamily\bfseries\color{night}#1}\\
  {\sffamily\footnotesize\color{summit}#2 · \href{#3}{\faExternalLink*~lien}}\\
  {\small #4}\par\smallskip}

\renewcommand{\partiecourante}{Coupole \versioncoupole}
\newcommand{\capture}[2]{\par\needspace{10\baselineskip}\begin{center}\includegraphics[width=#2\linewidth]{#1}\end{center}}
\newcommand{\cmd}[1]{\texttt{#1}}

\begin{document}
\begin{titlepage}
\centering
\vspace*{2.5cm}
\includegraphics[width=4.2cm]{../../logo/coupole_512.png}\par\vspace{1.2cm}
{\sffamily\bfseries\color{night}\fontsize{38}{44}\selectfont Coupole\par}
\vspace{0.4cm}
{\sffamily\Large\color{sky}Reference manual\par}
\vspace{0.3cm}
{\sffamily\large\color{summit}Version \versioncoupole{} --- \datecoupole\par}
\vspace{1.4cm}
\begin{minipage}{0.78\linewidth}\centering\sffamily\color{ink}
Free toolbox for students of the « Explorer et Comprendre l'Univers » diploma (Observatoire de Paris).\\[4pt]
OHP image bank · image quality · spectra and series · sites and times · cosmology.
\end{minipage}
\vfill
{\sffamily\small\color{summit}Author: ARP273-ROSE · GNU GPL licence version 3 or later\\
Data: « OHP student observations » bank, Observatoire de Paris / PADC (Etalab 2.0 licence)\par}
\end{titlepage}

\tableofcontents

% =====================================================================
\section{Overview}
Coupole is a free application (GPL-3.0) written in Python 3.10+, running on Windows, macOS and Linux, with a
graphical interface (PyQt6) and a complete command line. Every displayed text exists in French and English; the
system language is detected and can be forced (\emph{Language} menu, Preferences, or \cmd{--lang fr|en}).

\begin{center}\small
\begin{tabularx}{\linewidth}{@{}lX@{}}
\toprule
\textbf{Module} & \textbf{Purpose}\\
\midrule
OHP image bank & Inventory, catalogue, download, astrometric check, fixed headers, conversion, stackable sets, anomalies, sky map.\\
Image quality & Optional. FWHM, ellipticity, background, noise, gradient, saturation, trails, sampling.\\
Spectra and series & Reading and plotting 1D data (radio spectra, light curves, tables).\\
Cosmology & From redshift to distances, ages, volume, distance modulus, scale (Planck 2018 and others).\\
Sites and times & Site database, OpenStreetMap map, UTC and local time.\\
My computer & Hardware diagnosis; minimal example module.\\
\bottomrule
\end{tabularx}
\end{center}

Principles: all long work runs off the GUI thread (the interface stays responsive, with progress and
cancellation); an interrupted run resumes where it stopped; nothing is deleted without a trace; no data is sent
without consent; optional components (ASTAP, SEP, CuPy) may be missing.

% =====================================================================
\section{Installation}
\subsection{All-inclusive packages}
The packages bundle Python and every library (PyQt6, numpy, astropy, zstandard, lz4, SEP): nothing else to
install. They update themselves. ASTAP stays optional and is not included.

\begin{center}\small
\begin{tabularx}{\linewidth}{@{}l>{\ttfamily\small}lX@{}}
\toprule
\textbf{System} & \textbf{\normalfont File} & \textbf{Installation}\\
\midrule
Windows 10/11 64-bit & Coupole-Setup-x.y.z.exe & Double-click. SmartScreen warning at first launch (unsigned installer): « More info », « Run anyway ». Installed in the user profile, no administrator rights.\\
macOS Apple Silicon & Coupole-x.y.z-macos-arm64.dmg & Drag Coupole into Applications. First launch: right click, « Open ».\\
macOS Intel & Coupole-x.y.z-macos-x86\_64.dmg & Same.\\
Ubuntu / Debian x86\_64, ARM64 & coupole\_x.y.z\_amd64.deb & \texttt{sudo apt install ./coupole\_x.y.z\_amd64.deb} (\texttt{\_arm64.deb} on ARM): apt installs the Qt system libraries; \texttt{coupole} command, menu entry. Ubuntu 22.04 or newer, Debian 12 or newer. Uninstall: \texttt{sudo apt remove coupole}.\\
Other Linux x86\_64 / ARM64 & Coupole-x.y.z-linux-*.tar.gz & Unpack, run \texttt{installer.sh} (copied into \texttt{\textasciitilde/.local/share/coupole}, menu entry).\\
\bottomrule
\end{tabularx}
\end{center}

\subsection{With the computer's Python}
The scripts in the source folder look for Python 3.10 or newer, create a virtual environment, install the
dependencies and create launchers. One question only: Enter to confirm (\cmd{--oui} / \cmd{-Oui} to skip it).
\begin{ccode}
sh install.sh                \# Linux, macOS\\
install.bat                  \# Windows (double-click; runs install.ps1)
\end{ccode}
Without Python, or without the \texttt{venv} module (Debian: \texttt{python3-venv} package), the script says so and
explains how to get it. Locations: Linux \texttt{\textasciitilde/.local/share/coupole/venv}, launchers in
\texttt{\textasciitilde/.local/bin}; macOS \texttt{\textasciitilde/Library/Application Support/Coupole/venv} and
\texttt{\textasciitilde/Applications/Coupole.app}; Windows \texttt{\%LOCALAPPDATA\%\textbackslash
Programs\textbackslash CoupolePython}, Desktop and Start menu shortcuts.

\subsection{pip and pipx}
\begin{ccode}
pipx install git+https://github.com/ARP273-ROSE/coupole\\
pip install .                 \# from the sources; pip install ".[qualite]" adds SEP\\
python -m coupole             \# interface; coupole --help: command line
\end{ccode}
Extras: \texttt{qualite} (SEP), \texttt{validation} (lxml, xisf: XSD validation and second reader), \texttt{gpu}
(CuPy), \texttt{test} (pytest).

\subsection{Updates, uninstallation, data}
At startup (can be disabled) the packages check the latest release on GitHub and offer to install it: only the
code archive (\texttt{coupole-app-x.y.z.zip}, about 1 MB) is downloaded and unpacked by the application, the
previous version being restored on failure. A pip/pipx installation shows the command to run. The \texttt{.deb}
package does not update through the archive: its files (\texttt{/opt/coupole}) belong to dpkg; Coupole announces the
new version, gives the address of the new \texttt{.deb} and the \texttt{apt} command to run (\emph{Download the
package} button). Settings, cache and reports live outside the installation folder and survive updates as well as
removal:

\begin{center}\small
\begin{tabular}{@{}ll@{}}
\toprule
Windows & \texttt{\%LOCALAPPDATA\%\textbackslash Coupole}\\
macOS & \texttt{\textasciitilde/Library/Application Support/Coupole} (cache: \texttt{\textasciitilde/Library/Caches/Coupole})\\
Linux & \texttt{\textasciitilde/.config/coupole} (cache: \texttt{\textasciitilde/.cache/coupole})\\
All & \texttt{COUPOLE\_HOME} variable: replaces both (portable installations, tests)\\
\bottomrule
\end{tabular}
\end{center}

% =====================================================================
\section{Getting started}
The window has a module bar on the left (Ctrl+1 to Ctrl+9) and the current module on the right. Every interactive
element has a tooltip. \emph{Help} menu: help on the screen shown (F1), this manual (Shift+F1), keyboard shortcuts,
report a problem, incident reports, check for updates, About (version, licence, detected configuration).

\begin{center}\small
\begin{tabular}{@{}ll@{}}
\toprule
F1 & help on the screen\\ Shift+F1 & PDF manual\\ Ctrl+1 … Ctrl+9 & module 1 to 9\\ Ctrl+, & Preferences\\
Ctrl+Shift+A & ASTAP assistant\\ Ctrl+Shift+D & dark / light theme\\ Ctrl+R & OHP bank: refresh the inventory\\ Ctrl+Q & quit\\
\bottomrule
\end{tabular}
\end{center}

\capture{img/reglages_en.png}{0.62}
Preferences (Ctrl+,): interface language; language of folder names, object names and FITS comments;
appearance (section~\ref{sec:apparence}); output folder; format; parallelism; maximum rate; economy mode; update check; check of the bank for new images at startup and its maximum frequency (section~\ref{sec:nouveautes}). The \emph{Sources} tab lists the service
addresses (chapter~\ref{chap:evolutif}). An unreadable settings file is set aside (suffix \texttt{.corrompu-…}) and
the defaults take over: the application always starts.

\subsection{Appearance}\label{sec:apparence}
Coupole has its own theme, independent of the computer's (a system in dark mode no longer mixes its colours with the
application's: Fusion style, complete palette, fixed font size, WCAG contrast $\geq$ 4.5 for text, fields, buttons,
tooltips and selection). The \textbf{dark} theme is the default; the light theme remains available. It is changed
through the \emph{View} > \emph{Appearance} menu (checkable \emph{Light} / \emph{Dark} entries, Ctrl+Shift+D to
toggle) or in the Preferences: the same setting, applied at once and saved; both stay in sync. A user who had chosen
the light theme keeps it: the saved setting takes precedence over the default. The screenshots of this manual are
taken in the default theme.

% =====================================================================
\section{OHP image bank module}
\subsection{The bank}
« OHP student observations » gathers the images taken by the diploma's students at the Observatoire de
Haute-Provence: 7,989 FITS images on 8 October 2026, 166 objects under 194 names, 60 nights from 2015 to 2026,
78.0 GB; 1.20 m telescope (T120, 7,419 images, 1024 pixels of 0.77″, 13.1′ field) and the 50 cm IRIS robotic
telescope (570 images, 2048 pixels of 0.68″ or 4096 pixels of 0.46″, 23′ and 31′ fields). The images are already
calibrated (dark, flat) and astrometrically solved: \textbf{they must not be recalibrated}.

\subsection{Catalogue}
\capture{img/catalogue_en.png}{1.0}
The inventory comes from the Observatory's public TAP service (\texttt{http://tap-ufe.obspm.fr/tap}, table
\texttt{ivoa.obscore}, one synchronous ADQL query), kept in a cache; a snapshot is shipped for offline work.
\emph{Refresh the inventory} (Ctrl+R) queries the service again and flags new images and names (« new since … »).
Database names are grouped into canonical objects; selecting one or several objects lists their images on the
right, filtered by night, filter and doubtful dates; this list is the selection to process. The estimate gives the
volume to download, the expected output volume and the free space. Times are shown in UTC and in site time (hover:
also the user's time).

\paragraph{What you already own.} According to the output folder (\texttt{\_traitement/etat.sqlite}, read again in
the background when opening, after each run and whenever the folder changes, never blocking), the image list tells
apart with a marker and a soft text colour (readable in both themes, contrast $\geq 4.5$): \textbf{owned} (converted
and sorted, green check), \textbf{duplicate left out} (same pixels as an image already converted, grey dot),
\textbf{failed} (to retry, orange triangle) and \textbf{to download} (normal text, arrow); the marker's tooltip gives
the status, the date and the local path of the file. In the object table, the \emph{owned} column (« 120 / 300 »,
mini progress bar, sortable) and its marker (full disc: everything owned; half disc: partly; empty circle: nothing)
sum up each object. The \emph{To download only} box hides fully owned objects and, in the list, the images already
processed; the estimate then counts only what is missing (« 180 missing images out of 300: 1.5 GB »). A compact
legend under the list explains the markers. Command line: \cmd{coupole ohp inventory --missing [--dest FOLDER]
[--json]}, and \texttt{possedee}, \texttt{statut\_local}, \texttt{fichier\_local} columns in the CSV of \cmd{coupole
ohp list}.

\subsection{Processing}
\capture{img/traitement_en.png}{1.0}
For each image, duplicates excluded:
\begin{enumerate}[leftmargin=*,itemsep=1pt]
\item download: at most 4 at once (3 by default), total rate capped (8 MB/s by default), User-Agent naming Coupole,
  resume by HTTP Range requests, checks of size (server and inventory), multiple of 2,880 bytes and \texttt{SIMPLE}
  card;
\item reading of the FITS cards as written and of the pixels; SHA-1 fingerprint of the pixels (exact duplicates);
\item astrometric solution check (section~\ref{sec:astrometrie});
\item fixed header (section~\ref{sec:entetes});
\item writing in the chosen format, reading back and pixel-by-pixel comparison;
\item at the end: sorting into stacks, \texttt{LOT.txt}, \texttt{INDEX\_LOTS.csv}, \texttt{\_traitement/journal.csv}.
\end{enumerate}
The state is kept in \texttt{\_traitement/etat.sqlite}: \emph{Stop} then \emph{Start} (or running the same command
again) resumes without redoing anything. A folder keeps the name language and format of its first run.
\emph{Keep duplicates} also processes what would be left out; \emph{Keep the original FITS} stores them in
\texttt{\_traitement/fits\_origine}.

The chain is a \emph{pipeline}: downloads (threads) and conversions (separate processes, one per available core)
overlap, the queue of waiting FITS is bounded ($2\times$ conversions + downloads; when free space at the destination
is tight it is reduced automatically — never below conversions + 1 — and the log says so, rather than refusing the
job), and the overall rate is that of
the network or of the chosen cap (measured: 24 images of 8 MB at 8 MB/s in 24.5 s, i.e. the duration of the download
alone; without a cap, 86 MB/s of FITS converted on 6 cores). \emph{Pause} suspends feeding the pipeline (images in
progress finish) and \emph{Resume} restores it; \emph{Stop} interrupts downloads and conversions without waiting
(processes are terminated, ASTAP included: by signal on Linux and macOS, by a \emph{job object} on Windows). At the end, a report sums up images, duration, converted, duplicates,
failures (with their reason), volume written and stacks.

On a network share (NAS over SMB), where every access is a round trip, the state is committed at most once per second
and \texttt{JOURNAL.txt} is written in batches (a power cut only redoes the last second); the final sorting rewrites
only the \texttt{LOT.txt} files that changed and no longer walks the whole tree: running a processing again for one
image no longer costs a minute of sorting.

\subsection{Download everything, log, reorganise}\label{sec:tout}
\emph{Download everything…} takes the whole bank (about 8,000 images, 78 GB of FITS, 30 GB as XISF). On first use
(or when no output folder is set), Coupole asks for the destination folder through a native dialog and proposes
\texttt{Documents/Coupole/OHP\_DU\_ECU}; then it shows, \textbf{before any confirmation}, the volume to download,
the expected output volume, the estimated duration at the rate cap (78 GB at 8 MB/s: 2~h~42) and the free space; it
clearly refuses when space is short. Processing is the same as for any selection: resume after a cut or a close
(per-image state in \texttt{\_traitement/etat.sqlite}, already downloaded FITS found again), pause, end report.
Command line: \cmd{coupole ohp all --yes}.

\emph{Open the log} shows \texttt{\_traitement/JOURNAL.txt}: one timestamped (UTC), bilingual line per event ---
session start and end (resume, pause), each image downloaded, converted or left out with its reason, each error, each
name conflict; rotation at 5 MB (\texttt{JOURNAL-1.txt}, \texttt{JOURNAL-2.txt}). The tabular details stay in
\texttt{journal.csv}.

\emph{Reorganise files…} sorts into the stack tree images already converted by Coupole that live elsewhere or in an
older layout: each file is recognised by its header (source address \texttt{OHP:Source:URL}, or \texttt{HISTORY}
card « converted from … »), attached to the inventory and \textbf{moved} (never copied); a name already taken by a
foreign file gets a suffix and a log line; nothing is ever overwritten. Headers are read by several processes,
progress is shown (« 1,200 / 7,625 headers read ») and \emph{Stop} interrupts the reading (whatever is already
recognised is sorted). Command line: \cmd{coupole ohp reorganise FOLDER --dest DEST}.

\subsection{New images in the bank}\label{sec:nouveautes}
At startup (at most once a day, adjustable, can be disabled in Preferences), when the Internet is available and a
local copy exists in the output folder, Coupole compares a fresh TAP inventory with the copy. If there are new images,
a banner appears in the OHP bank (and a message in the status bar): « N new images (M objects, X GB) since DD/MM ---
Download now? / Later / Show ». \emph{Download} processes only the new images into the same tree; \emph{Show}
filters the catalogue; nothing is ever downloaded without explicit consent. The \emph{Check for new images} button
runs the comparison on demand. Command line: \cmd{coupole ohp new} (list) and \cmd{coupole ohp new --download}.

\subsection{Astrometric solution}\label{sec:astrometrie}
Consistency checks, always made: solution present and readable; scale within 2~\% of the nominal one
(\texttt{XPIXSZ}/\texttt{FOCALLEN}) or of the night median; angle within 5° of the night median (modulo 180°);
centre within 1.5° of the object medoid (fixed) or 2° of the object-and-night medoid (moving). Medians and medoids
come from the whole inventory.

With ASTAP (optional), each image is also solved independently:
\begin{center}\small
\begin{tabularx}{\linewidth}{@{}lX@{}}
\toprule
\textbf{Status} & \textbf{Meaning}\\
\midrule
confirmed & ASTAP finds the same centre within 10″ and the same scale within 1~\%.\\
validated & solution kept: checks pass, ASTAP silent (poor or crowded field, short exposure) or absent.\\
redone & missing or inconsistent solution replaced by ASTAP's (with SIP); the old one goes to HISTORY.\\
doubtful & checks fail and no ASTAP to decide: image stored apart.\\
failed & nothing usable: image stored in \texttt{\_sans\_solution\_astrometrique} (\texttt{\_no\_astrometric\_solution} with English names).\\
\bottomrule
\end{tabularx}
\end{center}

\subsection{Fixed headers}\label{sec:entetes}
All original cards are kept in order; every replaced value is copied into a \texttt{HISTORY original value: …}
card (French wording when names are in French).
\begin{center}\small
\begin{tabularx}{\linewidth}{@{}>{\ttfamily}lX@{}}
\toprule
\normalfont\textbf{Keyword} & \textbf{Fix}\\
\midrule
OBJECT & canonical name (database aliases merged).\\
FILTER, FILTSYS, FILTORIG & normalised filter (B, V, R, Ha, OIII, g, r, i, z), system (Johnson, Cousins, SDSS, Gunn, Narrowband), original value.\\
LATITUDE, LONGITUD & fixed \textbf{only when the swap is proven} (swapped values matching the site); correct: untouched; ambiguous: flagged, untouched. Idempotent.\\
SITELAT, SITELONG, SITEELEV & added when missing (MPC 511 site converted to WGS 84); present and consistent: kept; inconsistent: flagged.\\
RA, DEC, OBJCTRA, OBJCTDEC & true centre of the solution.\\
PIXSCALE, ORIENTAT & measured scale, position angle of the Y axis.\\
FOCALLEN, XPIXSZ & recomputed only when missing or inconsistent by more than 2~\%.\\
CTYPE1/2 & « -SIP » without coefficients (IRIS, PinPoint) brought back to TAN.\\
DATE-OBS, DATEDOUT & added when missing; \texttt{DATEDOUT = T} when the date is shared with another exposure.\\
XISFCONV & description of the conversion.\\
\bottomrule
\end{tabularx}
\end{center}
The \texttt{t\_min} column of the database equals \texttt{DATE-OBS} − exposure/2; \texttt{DATE-OBS} (exposure
start) prevails. A header date written in local time (whole-hour offset equal to the site time-zone offset against
\texttt{MJD-OBS}/\texttt{JD}, or a « night » exposure dated in daylight) is flagged in the log, never silently
fixed.

\subsection{Stacks}
\capture{img/lots_en.png}{0.9}
A stack gathers what is stacked together: same object, same instrument (T120, T120-bin1, T120-sCMOS, IRIS-2048,
IRIS-4096), same filter, and same field for a fixed object (all nights) or same night for a moving object.
\begin{ccode}
07\_Nebulae/NGC\_6888/field\_1\_T120/V/20250716-214427\_NGC\_6888\_V\_1s.xisf\\
01\_Asteroids/914\_Palisana/night\_2023-08-15\_T120/R/LOT.txt
\end{ccode}
\texttt{LOT.txt} (French then English) gives object, database names, instrument, filter, exposures and total
exposure, nights, mean centre, angle, solution status, doubtful dates and stacking advice (StarAlignment or
CometAlignment, rejection suited to the number of exposures). File names: UTC date and time of the exposure start,
object, filter, exposure; sorting by name = sorting in time. Names are safe under Windows. The \emph{Stacks} tab adds a \emph{complete} column: converted exposures against the
bank's exposures for the same key (object, instrument, filter, and night for a moving object); a fixed object split
into several fields counts its stacks together.

\subsection{Anomalies}
\capture{img/anomalies_en.png}{1.0}
\begin{center}\small
\begin{tabularx}{\linewidth}{@{}>{\ttfamily}llX@{}}
\toprule
\normalfont\textbf{Kind} & \textbf{Action} & \textbf{Meaning}\\
\midrule
meme\_fichier & left out & same exposure under another path or spelling\\
copie\_diurne & left out & daytime-dated copy of a correctly dated image\\
pixels\_identiques & left out & same pixels (SHA-1) as another image, at processing\\
memes\_meta & flagged & same date, filter, exposure, size, different name\\
date\_partagee & flagged & DATE-OBS identical to another exposure\\
date\_diurne & flagged & Sun above the horizon at the site (never for the Sun or radio data)\\
champ\_incoherent & flagged & centre far from the other exposures of the object\\
classement & flagged & name classified automatically or not at all\\
instrument & flagged & instrument missing from the table\\
\bottomrule
\end{tabularx}
\end{center}
On 8 October 2026: 133 + 11 = 144 duplicates left out, 1,654 shared dates, 3 daytime exposures, 92 inconsistent
fields. CSV export with bilingual columns; nothing is deleted.

\subsection{Sky map}
\capture{img/ciel_en.png}{0.95}
Aitoff projection, right ascension increasing to the left, 12~h at the centre; one point per fixed object (size:
number of images) and per night for moving objects; ecliptic dashed, galactic plane dotted.

\subsection{Online record (optional)}\label{sec:fiche}
\capture{img/fiche_en.png}{0.95}
When the Internet is available, the \emph{Online record} tab (or the button of the same name under the catalogue)
shows public information on the selected object:
\begin{itemize}
\item deep sky and stars: \textbf{SIMBAD} (CDS, Strasbourg) --- type, ICRS coordinates, magnitudes, parallax (and
the distance it gives when precise enough), preferred measured distance with its method and reference, radial
velocity and redshift, angular size, spectral and morphological types, other names; links to the SIMBAD page,
\textbf{Aladin Lite} (sky image centred on the object) and, for galaxies, \textbf{NED}; when SIMBAD does not know
the name, \textbf{Sesame} (CDS) also queries NED and VizieR;
\item asteroids, comets, trans-Neptunian objects: \textbf{JPL} Small-Body Database --- orbit class, near-Earth or
not, elements (a, e, i, q, Q, $\Omega$, $\omega$, period, epoch), diameter, albedo, rotation period, absolute
magnitude;
\item planets and NEOCP designations: no record (explanatory message).
\end{itemize}
\textbf{The identifier chosen by the service is always shown}: SIMBAD, for instance, attaches « NGC 6888 » to its
Wolf-Rayet star HD 192163; the record says so. A positive redshift can be sent to the Cosmology module
(\emph{Send this redshift to the Cosmology module} button).

\paragraph{Politeness and offline use.} Nothing is sent while the tab is hidden; one request per selected object
(after 0.6~s without another selection change), at most one per second and per service, 8~s timeout, Coupole's
User-Agent. Answers are kept for 30~days in Coupole's cache (\texttt{enligne.json}, with their date; « not found »
for one day). Offline, the last known record is shown with its date, otherwise a clear message; the rest of the
application never depends on it. The \emph{Query online services} box stops all requests. The addresses
(\texttt{simbad.base}, \texttt{simbad.page}, \texttt{aladin.page}, \texttt{ned.page}, \texttt{sbdb.page},
\texttt{resolution.sbdb}, \texttt{resolution.sesame}) are set like the other sources.

% =====================================================================
\section{ASTAP (optional)}
Without ASTAP, Coupole does everything else. ASTAP (Han Kleijn, MPL~2.0 licence,
\url{https://www.hnsky.org/astap.htm}) adds an independent check of each solution. It needs the program and one
star catalogue.

\capture{img/astap_en.png}{0.72}

\paragraph{Catalogue.} The official site states that for a field of 0.6° or more, D05, D20, D50 and D80 all suit,
and that dense catalogues serve small fields. The bank's fields range from 10.7′ (0.18°) to 31.4′ (0.52°):
\textbf{D80} is recommended (about 1.2 GB; catalogue of the reference processing, 6,888 solutions confirmed out of
7,625); D50 (about 0.9 GB) also suits IRIS.

\begin{center}\small
\begin{tabularx}{\linewidth}{@{}lX@{}}
\toprule
\textbf{System} & \textbf{Program (links of the official page)}\\
\midrule
Windows x64 & \texttt{astap\_setup.exe} then \texttt{d80\_star\_database.exe} (same folder, \texttt{C:\textbackslash Program Files\textbackslash astap}); optional: \texttt{astap\_command-line\_version\_win64.zip}\\
Windows ARM64 & \texttt{astap\_command-line\_version\_win11\_aarch64.zip} and \texttt{d80\_star\_database.zip}, in one folder\\
macOS Apple Silicon & \texttt{astap\_M1.pkg}, \texttt{d80\_star\_database.pkg}; once: \texttt{codesign --force -s - /Applications/ASTAP.app/Contents/MacOS/astap}\\
macOS Intel & \texttt{astap.pkg}, \texttt{d80\_star\_database.pkg}\\
Linux Debian/Ubuntu & \texttt{astap\_amd64.deb} (or \texttt{astap\_aarch64.deb}), \texttt{d80\_star\_database.deb} (into \texttt{/opt/astap})\\
Linux Fedora/openSUSE & \texttt{astap\_amd64.rpm}, catalogue \texttt{d80\_star\_database.zip} unzipped into \texttt{/opt/astap}\\
Linux Arch & \texttt{astap\_amd64.pkg.tar.zst}, catalogue as zip\\
Other Linux & \texttt{astap\_amd64.tar.gz} or \texttt{astap\_command-line\_version\_Linux\_amd64.zip}; no official AppImage\\
\bottomrule
\end{tabularx}
\end{center}

\begin{sloppypar}\noindent\textbf{Detection.} Order: chosen path (Preferences, \cmd{coupole astap --set}),
\texttt{COUPOLE\_ASTAP} variable, PATH (\texttt{astap\_cli}, then \texttt{astap}), default locations (Windows
\texttt{Program Files\textbackslash astap}, \texttt{\%LOCALAPPDATA\%\textbackslash astap}; macOS
\texttt{/Applications/ASTAP.app}, \texttt{/usr/local/opt/astap}; Linux \texttt{/opt/astap}, \texttt{/usr/bin},
\texttt{/usr/local/bin}, \texttt{\textasciitilde/astap}). Catalogue: chosen folder, \texttt{COUPOLE\_ASTAP\_CATALOGUE},
program folder, \texttt{/opt/astap}, \texttt{/usr/share/astap/data}, \texttt{/usr/local/opt/astap}. A catalogue is
declared complete with its 1,476 tiles; the version is read from \texttt{astap\_cli -h}.\end{sloppypar}

% =====================================================================
\section{Output formats}
\begin{center}\small
\begin{tabularx}{\linewidth}{@{}lXl@{}}
\toprule
\textbf{Format} & \textbf{Content} & \textbf{Readers}\\
\midrule
XISF (default) & Float32 in ADU, \texttt{bounds="-1000:65535"}; UInt16 for integer IRIS FITS; zstd+sh level 9 compression; \texttt{Observation:*}, \texttt{Instrument:*} properties & PixInsight\\
FITS \texttt{.fits.fz} & lossless tile compression: GZIP\_2 without quantization (floats), RICE\_1 (integers) & Siril, astropy, DS9\\
FITS & uncompressed float32 (or UInt16) & all\\
\bottomrule
\end{tabularx}
\end{center}
\textbf{XISF}: \texttt{bounds} sets the range PixInsight maps to [0,1] when opening: nothing to set, same scale for
every exposure; the −1,000 lower bound keeps the negative values of short calibrated exposures. Size: 36~\% of the
original FITS for the T120 (float64), 62~\% for IRIS. \textbf{Compressed FITS}: measured with astropy 8.0, RICE\_1
and HCOMPRESS\_1 quantize floats (one-ADU error, even with \texttt{quantize\_level=0}); only GZIP\_1 and GZIP\_2
without quantization restore floats bit for bit. \textbf{float32 FITS}: in PixInsight, set the FITS input range to
0–65,535 with \emph{Truncate}. For every format the difference to the original FITS stays below half a float32 ULP
(checked pixel by pixel).

% =====================================================================
\section{Image quality (optional)}
\capture{img/module_qualite_en.png}{1.0}
The module needs the SEP library (Source Extractor in Python, LGPL, about 1 MB); it never runs automatically
(\emph{Measure} button, \emph{Check image quality} box of the processing, or \cmd{coupole qualite FOLDER --write}).
Report per stack: \texttt{QUALITE.csv} and \texttt{QUALITE.txt}.

\paragraph{Speed, resume, sample.} Measurements run in separate processes (count given by the machine plan: cores,
memory, economy mode and the Preferences limit; at most 3 when the folder is on a network share, detected and
announced, since reading then limits the speed), one image being read while the others are computed; the interface
stays free, \emph{Stop} is immediate. Each measured image is kept at once in the cache
\texttt{\_traitement/qualite.sqlite} by (path, size, date) and shown; its stack's \texttt{QUALITE.csv} (atomic write)
is rewritten at most every 5~s and at the end of the stack or on stop: starting again redoes nothing, closing then
reopening resumes. The folder is walked only once (folders and dates read in parallel, cache read in one query):
starting again on the whole, already measured bank over a share takes 2~s instead of 72~s. Above 200 images, Coupole offers a \textbf{sample} of
5 images per stack (spread in time: first, last, middle…, adjustable in the tab and with \cmd{--sample N}), enough to
judge a stack, or \emph{Measure everything} with the duration estimated from the first 3 images; progress shows the
time left and the rate (images/min). On the author's NAS (Ryzen 5 8500G): 1.21 → 0.33~s per T120 image (vectorised
Moffat model, analytic derivatives, 120 stars per image), i.e. the whole bank in about ten minutes on 5 processes.

\begin{center}\small
\begin{tabularx}{\linewidth}{@{}lXl@{}}
\toprule
\textbf{Measure} & \textbf{Method} & \textbf{Synthetic validation}\\
\midrule
FWHM & elliptical Moffat with free $\beta$ fitted to each isolated, unsaturated star (model integrated over the pixel); median & ≤ 0.4~\% (Gauss and Moffat, 2.5–8 px)\\
ellipticity & $1-b/a$ of the same fit & ± 0.02\\
3 × 3 map & medians per ninth of the field & order recovered\\
background, noise & SEP, 64 px mesh & 0.5~\%; 5~\%\\
gradient, residual & plane fitted to the background map; residual amplitude & gradient and vignetting recovered\\
saturated & peak above SATURATE/DATAMAX (else 60,000 ADU) & 5 out of 5\\
trails & objects with major axis ≥ 12 px and $a/b \ge 6$ & 1 out of 1, 0 false\\
SNR & only with the gain (EGAIN/GAIN) & given or absent\\
sampling & FWHM < 2 px, 2–5 px, > 5 px & —\\
\bottomrule
\end{tabularx}
\end{center}
On 14 IRIS images (12 nights, FWHM from 2.4 to 11.3 px), the measured FWHM and the one written by PinPoint in the
header correlate (r = 0.974), median ratio 1.046 (0.85 to 1.16): the two estimators differ (PinPoint's method is
not described). Not offered for lack of validation: mag/arcsec² (no zero point), dust-donut detection. A FWHM
growing towards an edge may come from the sensor, field curvature, collimation, but also from seeing, tracking or
focus: the map alone cannot tell.

% =====================================================================
\section{Spectra and series}
\capture{img/module_donnees_en.png}{1.0}
Formats read: FITS (1D image with a linear \texttt{CRVAL/CDELT/CRPIX} WCS axis, 1 × 1 × N radio cubes, binary
tables) and CSV or text in columns (optional header, unit in parentheses or brackets). Radio velocity:
$v = c\,(1 - f/f_0)$ (FITS radio convention), $f_0$ read from \texttt{RESTFRQ}, otherwise the H\,I line at
1,420.405751768 MHz. The frame is that of the frequency axis (\texttt{SPECSYS}); no LSR or barycentric correction is
made implicitly. Formats of teaching radio antennas are added with a reader (CONTRIBUTING.md, « Adding a format »).

% =====================================================================
\section{Sites and times}
\capture{img/module_sites_en.png}{0.95}
Time is always kept in UTC (aware dates); the site's local time and the user's are shown as well. The « evening
date » of a night is the site's local date at the previous noon; daylight is computed from the Sun's altitude at the
site (astropy, no IERS download). Shipped sites: OHP (MPC 511), Meudon (005), Paris (007), geodetic coordinates
computed from the Minor Planet Center parallax constants. Partner sites are added (IANA time zone, MPC code if any)
once their coordinates are known from a reliable source. Map: OpenStreetMap tiles cached on disk, at most two
downloads at a time, User-Agent naming the application, attribution shown; offline, plain background and points. On
a dense screen (scaling $\geq 150$~\%, Mac Retina) the map draws the tiles of the next zoom level at half size (four
per logical tile, the « @2x » equivalent): it stays sharp instead of being upscaled by the system.

% =====================================================================
\section{Cosmology}\label{sec:cosmo}
\capture{img/module_cosmo_en.png}{1.0}
The redshift $z$ of a distant object does not give « one » distance but several, all correct, answering different
questions. The module computes them in the $\Lambda$CDM model with
\[ E(z)^2 = \Omega_r(z)(1+z)^4 + \Omega_m(1+z)^3 + \Omega_k(1+z)^2 + \Omega_\Lambda, \qquad
   D_C = \frac{c}{H_0}\int_0^z\frac{dz'}{E(z')}, \]
radiation included (CMB photons and neutrinos, one of them massive, 0.06~eV).
\begin{center}
\begin{tabularx}{\linewidth}{@{}lX@{}}
\toprule
Quantity & Use \\
\midrule
$D_C$ line-of-sight comoving & present distance of the object; saturates at $\approx 46$~Gly (particle horizon)\\
$D_M$ transverse comoving & equals $D_C$ in a flat universe; enters $D_L$ and $D_A$\\
$D_L = (1+z)D_M$ & photometry: $F = L/(4\pi D_L^2)$; distance modulus $\mu = 5\log_{10}(D_L/10~\mathrm{pc})$\\
$D_A = D_M/(1+z)$ & sizes: $\theta = \ell/D_A$; scale in kpc per arcsecond; maximum near $z\approx1.6$\\
$c\,t_L$ & path travelled by the light (the most intuitive, the least physical)\\
$t_L$, age at $z$, $t_0$ & lookback time; check $t_L + t(z) = t_0$\\
$V_C$ & comoving volume out to $z$ (Hogg 1999, exact in curved universes too)\\
velocities & $cz$ (naive), relativistic Doppler (inappropriate), $H_0 D_C$ (the only correct one)\\
\bottomrule
\end{tabularx}
\end{center}

\paragraph{Parameters.} Planck 2018 by default ($H_0 = 67.66$, $\Omega_m = 0.3111$ incl. $\Omega_\nu = 0.00144$,
$T_0 = 2.7255$~K) with the $1\sigma$ uncertainty propagated from \emph{correlated} $H_0$ and $\Omega_m$
($\rho = -0.976$: ignoring it would overestimate the uncertainty) and, optionally, the comparison with SH0ES
($H_0 = 73.04$, Hubble tension). For comparison: Planck 2015, WMAP 9-year, a « textbook » $\Lambda$CDM
(70~/ 0.3~/ 0.7 without radiation) and a custom set ($H_0$, $\Omega_m$, $\Omega_k$, Planck 2018 radiation).
Ranges: $H_0$ from 40 to 100, $\Omega_m$ from 0.1 to 1, $\Omega_k$ from $-0.3$ to $0.3$ ($\pm 0.05$ for Planck
2018); an out-of-range value is \emph{refused with an explanation}, never silently corrected.

\paragraph{Input.} A $z$ (comma or point), one of the examples (M\,87, 3C\,273, $z = 1$, $z = 2.34$, ULAS J1120,
GN-z11, reionisation, CMB), or an object name: its redshift is requested from SIMBAD (tolerant naming: « m87 »,
« ngc4486 »; several candidates: you choose). $z \le 0$ (approaching object, like M\,31) gives no distance; below
$z = 0.03$ the object's peculiar velocity dominates (warning); beyond $z \approx 1090$ the universe is opaque. Table
and curves export to CSV (« ; », UTF-8; raw values in Mpc, Gyr, Gpc$^3$).

\paragraph{Slider.} Under the $z$ field, a slider spans the whole range on a logarithmic scale (0.001 to 1100,
marks at 0.01 / 0.1 / 1 / 10 / 100 / 1000). While dragging, the vertical marker of the curves and the table of values
follow live: values are then \emph{interpolated} on the curve grid (300 logarithmic points, typical deviation
$< 0.5$~\%, $\sigma$ and SH0ES columns suspended) and the status line says so; on release the exact computation
resumes. Arrows: fine step (1~\% of a decade); Page up/down: one decade. The $z$ field and the slider stay in sync
both ways. The curves follow the window when it is resized, including when they sit under the table (window
narrower than 1500 px): half of the visible height goes to them.

\paragraph{Origin and checks.} The computation core is taken from the same author's « cosmologie-redshift »
calculator (public repository \url{https://github.com/ARP273-ROSE/cosmologie-redshift}), with astropy.cosmology.
Before integration it was compared on 8 October 2026 with astropy and with an independent SageMath integration
(mpmath, 18 digits, exact Fermi-Dirac function for the neutrinos), from $z = 10^{-8}$ to 1100, for
$\Omega_k = 0, \pm0.01, \pm0.05$: relative deviation $\le 3\times10^{-6}$ on distances and $\le 2\times10^{-5}$ on
ages and $E(z)$ (Komatsu fit of the neutrinos in astropy); no bug. astropy's comoving volume, however, loses all
its digits at very small $z$ in curved universes (down to a negative volume): Coupole uses a series expansion there,
matched to the closed form. Scripts: \texttt{outils/verif\_cosmo/}.

\begin{ccode}
coupole cosmo 0.158 --shoes\\
coupole cosmo --object "3C 273"\\
coupole cosmo 1 2 3 --model wmap9 --csv results.csv\\
coupole cosmo --curves 0.001 1500 300 --csv curves.csv
\end{ccode}

% =====================================================================
\section{Computer and performance}
\capture{img/module_machine_en.png}{0.9}
Detection without mandatory dependency: system, architecture, processor, physical and logical cores (bounded by
affinity and container quota), total and available memory, free space, graphics cards (nvidia-smi,
\texttt{/sys/class/drm}, system\_profiler, CIM). Parallelism: 3 downloads (at most 4), conversions = physical
cores − 1 bounded by available memory at 500 MB per conversion (peak measured on an IRIS 4096² image, at most half
of the free memory); economy mode (1 conversion, 2 downloads) automatic below 4 GB or 2 cores; manual limits, also
bounded by memory. Examples: 2 cores and 2 GB → 2 downloads, 1 conversion; 6 cores and 20 GB free → 3 and 5;
32 cores and 60 GB free → 3 and 16 (cap). The graphics card is not used by the OHP bank (work bound by network, disk
and Zstandard); \texttt{coupole.core.calcul} offers CuPy to future modules.

Nothing heavy runs on the GUI thread: astropy, scipy and SEP are loaded only when first needed, the machine probes and
the ASTAP detection run in the background, and closing the window stops then waits for every thread, task and process
(see \texttt{docs/AUDIT\_2026-10.md}). Large lists stay smooth at 80,000 rows: cells are computed only for the rows
shown, sorting uses one key per row (not pairwise comparisons), selecting no longer makes every drawing walk all rows,
and changing the theme only recomputes colours (see \texttt{docs/AUDIT2\_2026-10.md}).

\begin{center}\small
\begin{tabular}{@{}lr@{}}
\toprule
\textbf{Measurement (Ryzen 5 8500G, 6 cores, NVMe; Debian container)} & \textbf{Value}\\
\midrule
\texttt{coupole --version} & 0.1–0.2 s\\
window shown (astropy, scipy, SEP not loaded) & 0.42 s (1.5 s before the audit)\\
inventory loaded (background thread) & 3.3 s at first launch, 0.8 s afterwards\\
full chain, 24 T120 images of 8 MB, 8 MB/s cap & 24.5 s = duration of the download alone (24.2 s)\\
same chain without a cap (local server) & 2.3 s, 86 MB/s of FITS, 10 images/s\\
conversion of one T120 image (one process) & 0.13 s (zstd 14 ms, reading 18 ms, writing and reading back 29 ms)\\
memory peak of a T120 / IRIS 4096² conversion (dedicated process) & 163 MB / 365 MB (170 / 573 before)\\
interface during 3 downloads (28 MB/s) & max gap of a 20 ms timer: 20.8 ms\\
\midrule
\multicolumn{2}{@{}l}{\emph{Second audit (0.1.4): real bank of 7,625 XISF; network share simulated at 2 ms per access}}\\
final sorting with nothing to move, on the share & 1.5 s (60.7 s before)\\
48 images to a destination on the share & 4.9 s (20.3 s before)\\
quality, starting again on the already measured bank (share) & 1.9 s (72 s before)\\
reorganise the bank (reading the 7,625 headers) & 57 s (243 s before)\\
select all, 80,000 images (10 × the bank) & 0.2 s (3.1 s before; 0.5 s → 0.17 s on the bank)\\
sort 80,000 images & 0.1–0.2 s (3.7–4.1 s before)\\
\bottomrule
\end{tabular}
\end{center}

% =====================================================================
\section{Command line}
Everything the interface does is available on the command line; \cmd{--json} for scripts. Without argument,
\cmd{coupole} opens the interface. Commands have a French name and an English alias.
\begin{ccode}
coupole --lang en ohp estimate "NGC 6888" --night 2025-07-16\\
coupole --lang en ohp process "NGC 6888" --night 2025-07-16 --filter B --filter V --filter R --dest \textasciitilde/OHP --yes\\
coupole --lang en ohp download Pluto --no-doubtful-dates --dest \textasciitilde/FITS\\
coupole --lang en qualite \textasciitilde/OHP --write
\end{ccode}
{\small% Généré par outils/tables_cli.py — ne pas modifier à la main.
\begin{tabularx}{\linewidth}{@{}>{\ttfamily\small\raggedright}p{0.38\linewidth}X@{}}
\toprule
\multicolumn{2}{@{}p{\linewidth}@{}}{\sffamily\bfseries coupole}\\
\midrule
--version & show the version and exit\\
--lang,\allowbreak{} --langue \{auto,\allowbreak{}fr,\allowbreak{}en\} & message language (auto: system language)\\
-v,\allowbreak{} --verbeux,\allowbreak{} --verbose & detailed log\\
--reinitialiser-interface,\allowbreak{} --reset-interface & erase the saved layout (window, columns, filters, file-dialog folders) and exit; settings do not change\\
\bottomrule
\end{tabularx}
\par\medskip
\begin{tabularx}{\linewidth}{@{}>{\ttfamily\small\raggedright}p{0.38\linewidth}X@{}}
\toprule
\multicolumn{2}{@{}p{\linewidth}@{}}{\sffamily\bfseries coupole astap\newline\normalfont\small Look for ASTAP and its star catalogue, show what to install for this system, or record their location. Coupole works without ASTAP.}\\
\midrule
--guide & show the installation guide for this system\\
--definir,\allowbreak{} --set PATH & path of the ASTAP executable (or its folder)\\
--catalogue,\allowbreak{} --catalog FOLDER & star catalogue folder (D80, D50...)\\
--oublier,\allowbreak{} --forget & forget recorded paths (back to detection)\\
--json & JSON output (for scripts)\\
\bottomrule
\end{tabularx}
\par\medskip
\begin{tabularx}{\linewidth}{@{}>{\ttfamily\small\raggedright}p{0.38\linewidth}X@{}}
\toprule
\multicolumn{2}{@{}p{\linewidth}@{}}{\sffamily\bfseries coupole rapports (reports)\newline\normalfont\small Nothing is sent without consent. Without consent, reports stay in a local folder.}\\
\midrule
--activer,\allowbreak{} --enable & allow sending (anonymous) reports\\
--desactiver,\allowbreak{} --disable & forbid sending reports\\
\bottomrule
\end{tabularx}
\par\medskip
\begin{tabularx}{\linewidth}{@{}>{\ttfamily\small\raggedright}p{0.38\linewidth}X@{}}
\toprule
\multicolumn{2}{@{}p{\linewidth}@{}}{\sffamily\bfseries coupole maj (update)\newline\normalfont\small Query the releases published on GitHub. The standalone package updates itself; a pip/pipx installation shows the command to run.}\\
\midrule
--appliquer,\allowbreak{} --apply & download and install the update (standalone package)\\
\bottomrule
\end{tabularx}
\par\medskip
\begin{tabularx}{\linewidth}{@{}>{\ttfamily\small\raggedright}p{0.38\linewidth}X@{}}
\toprule
\multicolumn{2}{@{}p{\linewidth}@{}}{\sffamily\bfseries coupole sources\newline\normalfont\small Without option: list the addresses and their origin (shipped default, remote file, forced value).}\\
\midrule
--forcer,\allowbreak{} --set KEY VALUE & force a value (never overwritten by an update)\\
--defaut,\allowbreak{} --default KEY & go back to the default value for this key\\
--tout-reinitialiser,\allowbreak{} --reset-all & go back to the default values for every key\\
--tester,\allowbreak{} --test KEY & test the connection to this address\\
--distant,\allowbreak{} --remote & fetch the sources file published in the repository now\\
\bottomrule
\end{tabularx}
\par\medskip
\begin{tabularx}{\linewidth}{@{}>{\ttfamily\small\raggedright}p{0.38\linewidth}X@{}}
\toprule
\multicolumn{2}{@{}p{\linewidth}@{}}{\sffamily\bfseries coupole ohp inventaire (inventory)\newline\normalfont\small Bank inventory from the public TAP service (http://tap-ufe.obspm.fr/tap, table ivoa.obscore), kept in a local cache.}\\
\midrule
--rafraichir,\allowbreak{} --refresh & query the Observatory's TAP service again (one request)\\
--manquantes,\allowbreak{} --missing & list what remains to download compared with the output folder (--dest)\\
--dest FOLDER & destination folder (default: \textasciitilde{}/Coupole/OHP\_DU\_ECU or the one in the settings)\\
--json & JSON output (for scripts)\\
\bottomrule
\end{tabularx}
\par\medskip
\begin{tabularx}{\linewidth}{@{}>{\ttfamily\small\raggedright}p{0.38\linewidth}X@{}}
\toprule
\multicolumn{2}{@{}p{\linewidth}@{}}{\sffamily\bfseries coupole ohp catalogue (catalog)\newline\normalfont\small The 194 database names grouped into objects, sorted by type.}\\
\midrule
--type TYPE & category: ast, neocp, com, pla, tno, pn, neb, amas, gal, eto (repeatable)\\
--telescope \{T120,\allowbreak{}IRIS\} & telescope (repeatable)\\
--chercher,\allowbreak{} --search TEXT & keep objects whose name contains this text\\
--json & JSON output (for scripts)\\
\bottomrule
\end{tabularx}
\par\medskip
\begin{tabularx}{\linewidth}{@{}>{\ttfamily\small\raggedright}p{0.38\linewidth}X@{}}
\toprule
\multicolumn{2}{@{}p{\linewidth}@{}}{\sffamily\bfseries coupole ohp images (list)\newline\normalfont\small One line per image: date, telescope, filter, exposure, flags (duplicate, shared date, daytime).}\\
\midrule
OBJECT & object(s): « NGC 6888 », « Pluto », « 1998 AG6 », or a database name\\
--type TYPE & category: ast, neocp, com, pla, tno, pn, neb, amas, gal, eto (repeatable)\\
--telescope \{T120,\allowbreak{}IRIS\} & telescope (repeatable)\\
--filtre,\allowbreak{} --filter F & filter: B, V, R (Johnson and Cousins), Bc, Vc, Rc, Ha, OIII, g, r, i, z, or the exact name (« Johnson R »); repeatable or comma-separated\\
--nuit,\allowbreak{} --night YYYY-MM-DD & night (evening date); repeatable\\
--sans-dates-douteuses,\allowbreak{} --no-doubtful-dates & leave out exposures with a shared DATE-OBS or dated in daytime\\
--tout,\allowbreak{} --all & the whole bank (≈ 78 GB)\\
--csv FILE & also write the list to this CSV file\\
--dest FOLDER & destination folder (default: \textasciitilde{}/Coupole/OHP\_DU\_ECU or the one in the settings)\\
\bottomrule
\end{tabularx}
\par\medskip
\begin{tabularx}{\linewidth}{@{}>{\ttfamily\small\raggedright}p{0.38\linewidth}X@{}}
\toprule
\multicolumn{2}{@{}p{\linewidth}@{}}{\sffamily\bfseries coupole ohp estimer (estimate)\newline\normalfont\small Estimate before downloading, with the free space at the destination.}\\
\midrule
OBJECT & object(s): « NGC 6888 », « Pluto », « 1998 AG6 », or a database name\\
--type TYPE & category: ast, neocp, com, pla, tno, pn, neb, amas, gal, eto (repeatable)\\
--telescope \{T120,\allowbreak{}IRIS\} & telescope (repeatable)\\
--filtre,\allowbreak{} --filter F & filter: B, V, R (Johnson and Cousins), Bc, Vc, Rc, Ha, OIII, g, r, i, z, or the exact name (« Johnson R »); repeatable or comma-separated\\
--nuit,\allowbreak{} --night YYYY-MM-DD & night (evening date); repeatable\\
--sans-dates-douteuses,\allowbreak{} --no-doubtful-dates & leave out exposures with a shared DATE-OBS or dated in daytime\\
--tout,\allowbreak{} --all & the whole bank (≈ 78 GB)\\
--format \{xisf,\allowbreak{}xisf16,\allowbreak{}fz,\allowbreak{}fits\} & output format: xisf (PixInsight, default), xisf16 (XISF compatible with N.I.N.A. and Siril: 16-bit integers, zlib), fz (lossless compressed FITS), fits (float32)\\
--dest FOLDER & destination folder (default: \textasciitilde{}/Coupole/OHP\_DU\_ECU or the one in the settings)\\
\bottomrule
\end{tabularx}
\par\medskip
\begin{tabularx}{\linewidth}{@{}>{\ttfamily\small\raggedright}p{0.38\linewidth}X@{}}
\toprule
\multicolumn{2}{@{}p{\linewidth}@{}}{\sffamily\bfseries coupole ohp traiter (process)\newline\normalfont\small Full processing of a selection, resumable: running the same command again resumes where it stopped.}\\
\midrule
OBJECT & object(s): « NGC 6888 », « Pluto », « 1998 AG6 », or a database name\\
--type TYPE & category: ast, neocp, com, pla, tno, pn, neb, amas, gal, eto (repeatable)\\
--telescope \{T120,\allowbreak{}IRIS\} & telescope (repeatable)\\
--filtre,\allowbreak{} --filter F & filter: B, V, R (Johnson and Cousins), Bc, Vc, Rc, Ha, OIII, g, r, i, z, or the exact name (« Johnson R »); repeatable or comma-separated\\
--nuit,\allowbreak{} --night YYYY-MM-DD & night (evening date); repeatable\\
--sans-dates-douteuses,\allowbreak{} --no-doubtful-dates & leave out exposures with a shared DATE-OBS or dated in daytime\\
--tout,\allowbreak{} --all & the whole bank (≈ 78 GB)\\
--dest FOLDER & destination folder (default: \textasciitilde{}/Coupole/OHP\_DU\_ECU or the one in the settings)\\
--format \{xisf,\allowbreak{}xisf16,\allowbreak{}fz,\allowbreak{}fits\} & output format: xisf (PixInsight, default), xisf16 (XISF compatible with N.I.N.A. and Siril: 16-bit integers, zlib), fz (lossless compressed FITS), fits (float32)\\
--noms,\allowbreak{} --names \{fr,\allowbreak{}en\} & language of folder names, object names and headers (default: interface language)\\
--astap \{auto,\allowbreak{}tous,\allowbreak{}suspectes,\allowbreak{}jamais,\allowbreak{}all,\allowbreak{}suspicious,\allowbreak{}never\} & ASTAP check: auto/all (every image when ASTAP is present), suspicious, never\\
--telechargements,\allowbreak{} --downloads N & simultaneous downloads (automatic by default; at most 4)\\
--conversions N & simultaneous conversions (default: from cores and memory)\\
--econome,\allowbreak{} --eco & economy mode: 1 conversion, 2 downloads\\
--debit,\allowbreak{} --rate MB/s & maximum download rate (default 8 MB/s)\\
--garder-fits,\allowbreak{} --keep-fits & also keep the original FITS (in \_traitement/fits\_origine)\\
--garder-doublons,\allowbreak{} --keep-duplicates & also process duplicates (inventory and identical pixels) instead of leaving them out\\
--oui,\allowbreak{} --yes & accept a selection over 5 GB without confirmation\\
\bottomrule
\end{tabularx}
\par\medskip
\begin{tabularx}{\linewidth}{@{}>{\ttfamily\small\raggedright}p{0.38\linewidth}X@{}}
\toprule
\multicolumn{2}{@{}p{\linewidth}@{}}{\sffamily\bfseries coupole ohp tout (all)\newline\normalfont\small The whole bank (≈ 8,000 images, ≈ 78 GB of FITS → ≈ 30 GB as XISF), with a volume and time estimate before confirmation; running the same command again resumes where it stopped.}\\
\midrule
--dest FOLDER & destination folder (default: \textasciitilde{}/Coupole/OHP\_DU\_ECU or the one in the settings)\\
--format \{xisf,\allowbreak{}xisf16,\allowbreak{}fz,\allowbreak{}fits\} & output format: xisf (PixInsight, default), xisf16 (XISF compatible with N.I.N.A. and Siril: 16-bit integers, zlib), fz (lossless compressed FITS), fits (float32)\\
--noms,\allowbreak{} --names \{fr,\allowbreak{}en\} & language of folder names, object names and headers (default: interface language)\\
--astap \{auto,\allowbreak{}tous,\allowbreak{}suspectes,\allowbreak{}jamais,\allowbreak{}all,\allowbreak{}suspicious,\allowbreak{}never\} & ASTAP check: auto/all (every image when ASTAP is present), suspicious, never\\
--telechargements,\allowbreak{} --downloads N & simultaneous downloads (automatic by default; at most 4)\\
--conversions N & simultaneous conversions (default: from cores and memory)\\
--econome,\allowbreak{} --eco & economy mode: 1 conversion, 2 downloads\\
--debit,\allowbreak{} --rate MB/s & maximum download rate (default 8 MB/s)\\
--garder-fits,\allowbreak{} --keep-fits & also keep the original FITS (in \_traitement/fits\_origine)\\
--garder-doublons,\allowbreak{} --keep-duplicates & also process duplicates (inventory and identical pixels) instead of leaving them out\\
--oui,\allowbreak{} --yes & accept a selection over 5 GB without confirmation\\
\bottomrule
\end{tabularx}
\par\medskip
\begin{tabularx}{\linewidth}{@{}>{\ttfamily\small\raggedright}p{0.38\linewidth}X@{}}
\toprule
\multicolumn{2}{@{}p{\linewidth}@{}}{\sffamily\bfseries coupole ohp nouveautes (new)\newline\normalfont\small Query the TAP service again, compare with the local copy (output folder) and list what is new; --download processes them into the same tree (never without this option).}\\
\midrule
--telecharger,\allowbreak{} --download & download and convert the listed new images\\
--dest FOLDER & destination folder (default: \textasciitilde{}/Coupole/OHP\_DU\_ECU or the one in the settings)\\
--format \{xisf,\allowbreak{}xisf16,\allowbreak{}fz,\allowbreak{}fits\} & output format: xisf (PixInsight, default), xisf16 (XISF compatible with N.I.N.A. and Siril: 16-bit integers, zlib), fz (lossless compressed FITS), fits (float32)\\
--noms,\allowbreak{} --names \{fr,\allowbreak{}en\} & language of folder names, object names and headers (default: interface language)\\
--astap \{auto,\allowbreak{}tous,\allowbreak{}suspectes,\allowbreak{}jamais,\allowbreak{}all,\allowbreak{}suspicious,\allowbreak{}never\} & ASTAP check: auto/all (every image when ASTAP is present), suspicious, never\\
--telechargements,\allowbreak{} --downloads N & simultaneous downloads (automatic by default; at most 4)\\
--conversions N & simultaneous conversions (default: from cores and memory)\\
--econome,\allowbreak{} --eco & economy mode: 1 conversion, 2 downloads\\
--debit,\allowbreak{} --rate MB/s & maximum download rate (default 8 MB/s)\\
--garder-fits,\allowbreak{} --keep-fits & also keep the original FITS (in \_traitement/fits\_origine)\\
--garder-doublons,\allowbreak{} --keep-duplicates & also process duplicates (inventory and identical pixels) instead of leaving them out\\
--oui,\allowbreak{} --yes & accept a selection over 5 GB without confirmation\\
\bottomrule
\end{tabularx}
\par\medskip
\begin{tabularx}{\linewidth}{@{}>{\ttfamily\small\raggedright}p{0.38\linewidth}X@{}}
\toprule
\multicolumn{2}{@{}p{\linewidth}@{}}{\sffamily\bfseries coupole ohp reorganiser (reorganise, reorganize)\newline\normalfont\small Walk a folder, recognise files produced by Coupole (header) and move them into the stack tree of --dest; name conflict → suffix and a log line; nothing is overwritten.}\\
\midrule
FOLDER & folder holding the files to sort\\
--dest FOLDER & destination folder (default: \textasciitilde{}/Coupole/OHP\_DU\_ECU or the one in the settings)\\
\bottomrule
\end{tabularx}
\par\medskip
\begin{tabularx}{\linewidth}{@{}>{\ttfamily\small\raggedright}p{0.38\linewidth}X@{}}
\toprule
\multicolumn{2}{@{}p{\linewidth}@{}}{\sffamily\bfseries coupole ohp telecharger (download)\newline\normalfont\small Layout: <folder>/<object>/<night>\_<telescope>/<original file>.}\\
\midrule
OBJECT & object(s): « NGC 6888 », « Pluto », « 1998 AG6 », or a database name\\
--type TYPE & category: ast, neocp, com, pla, tno, pn, neb, amas, gal, eto (repeatable)\\
--telescope \{T120,\allowbreak{}IRIS\} & telescope (repeatable)\\
--filtre,\allowbreak{} --filter F & filter: B, V, R (Johnson and Cousins), Bc, Vc, Rc, Ha, OIII, g, r, i, z, or the exact name (« Johnson R »); repeatable or comma-separated\\
--nuit,\allowbreak{} --night YYYY-MM-DD & night (evening date); repeatable\\
--sans-dates-douteuses,\allowbreak{} --no-doubtful-dates & leave out exposures with a shared DATE-OBS or dated in daytime\\
--tout,\allowbreak{} --all & the whole bank (≈ 78 GB)\\
--dest FOLDER & destination folder (default: \textasciitilde{}/Coupole/OHP\_DU\_ECU or the one in the settings)\\
--telechargements,\allowbreak{} --downloads N & simultaneous downloads (automatic by default; at most 4)\\
--debit,\allowbreak{} --rate MB/s & maximum download rate (default 8 MB/s)\\
--oui,\allowbreak{} --yes & accept a selection over 5 GB without confirmation\\
\bottomrule
\end{tabularx}
\par\medskip
\begin{tabularx}{\linewidth}{@{}>{\ttfamily\small\raggedright}p{0.38\linewidth}X@{}}
\toprule
\multicolumn{2}{@{}p{\linewidth}@{}}{\sffamily\bfseries coupole ohp anomalies (anomalies-report)\newline\normalfont\small Duplicates (same file, daytime copy, identical pixels), shared or daytime dates, inconsistent field, classification to check, unknown instrument. Nothing is deleted.}\\
\midrule
--dest FOLDER & add the pixel duplicates found in this processing folder\\
--csv FILE & also write the list to this CSV file\\
--nouveaux,\allowbreak{} --new & only images new since the last refresh\\
\bottomrule
\end{tabularx}
\par\medskip
\begin{tabularx}{\linewidth}{@{}>{\ttfamily\small\raggedright}p{0.38\linewidth}X@{}}
\toprule
\multicolumn{2}{@{}p{\linewidth}@{}}{\sffamily\bfseries coupole ohp classer (classify)\newline\normalfont\small Without argument: list the names to check. With a database name and --object: record the correction (merge into an existing object if --type is omitted). --online: query JPL SBDB and CDS Sesame.}\\
\midrule
DB\_NAME & target name as it appears in the database\\
--objet,\allowbreak{} --object OBJECT & canonical object to associate (existing one: merge)\\
--type \{ast,\allowbreak{}neocp,\allowbreak{}com,\allowbreak{}pla,\allowbreak{}tno,\allowbreak{}pn,\allowbreak{}neb,\allowbreak{}amas,\allowbreak{}gal,\allowbreak{}eto,\allowbreak{}autre\} & category: ast, neocp, com, pla, tno, pn, neb, amas, gal, eto (repeatable)\\
--oublier,\allowbreak{} --forget & forget the correction of this name\\
--en-ligne,\allowbreak{} --online & resolve unknown names online (JPL SBDB, CDS Sesame), with cache\\
\bottomrule
\end{tabularx}
\par\medskip
\begin{tabularx}{\linewidth}{@{}>{\ttfamily\small\raggedright}p{0.38\linewidth}X@{}}
\toprule
\multicolumn{2}{@{}p{\linewidth}@{}}{\sffamily\bfseries coupole ohp ranger (sort)\newline\normalfont\small Useful after an interruption, or to merge several runs in the same folder.}\\
\midrule
--dest FOLDER & destination folder (default: \textasciitilde{}/Coupole/OHP\_DU\_ECU or the one in the settings)\\
\bottomrule
\end{tabularx}
\par\medskip
\begin{tabularx}{\linewidth}{@{}>{\ttfamily\small\raggedright}p{0.38\linewidth}X@{}}
\toprule
\multicolumn{2}{@{}p{\linewidth}@{}}{\sffamily\bfseries coupole ohp bilan (status)\newline\normalfont\small progress of a processing folder}\\
\midrule
--dest FOLDER & destination folder (default: \textasciitilde{}/Coupole/OHP\_DU\_ECU or the one in the settings)\\
--json & JSON output (for scripts)\\
\bottomrule
\end{tabularx}
\par\medskip
\begin{tabularx}{\linewidth}{@{}>{\ttfamily\small\raggedright}p{0.38\linewidth}X@{}}
\toprule
\multicolumn{2}{@{}p{\linewidth}@{}}{\sffamily\bfseries coupole ohp fusionner (merge)\newline\normalfont\small merge the state database on the share with the local working database (after a divergence)}\\
\midrule
--dest FOLDER & destination folder (default: \textasciitilde{}/Coupole/OHP\_DU\_ECU or the one in the settings)\\
\bottomrule
\end{tabularx}
\par\medskip
\begin{tabularx}{\linewidth}{@{}>{\ttfamily\small\raggedright}p{0.38\linewidth}X@{}}
\toprule
\multicolumn{2}{@{}p{\linewidth}@{}}{\sffamily\bfseries coupole ohp metadonnees (metadata)\newline\normalfont\small For each XISF/FITS of the folder: FOCALLEN matched to the measured scale (old value in HISTORY) and the XISF properties Instrument:Sensor:XPixelSize, Instrument:Camera:XBinning, Instrument:Telescope:FocalLength… Pixels are not touched (read back and compared), atomic writing, log in \_traitement/metadonnees.csv. Without --rewrite: dry run.}\\
\midrule
FOLDER & output folder (or one stack)\\
--reecrire,\allowbreak{} --rewrite & really write (otherwise: only say what would change)\\
--noms,\allowbreak{} --names \{fr,\allowbreak{}en\} & language of folder names, object names and headers (default: interface language)\\
\bottomrule
\end{tabularx}
\par\medskip
\begin{tabularx}{\linewidth}{@{}>{\ttfamily\small\raggedright}p{0.38\linewidth}X@{}}
\toprule
\multicolumn{2}{@{}p{\linewidth}@{}}{\sffamily\bfseries coupole qualite\newline\normalfont\small Optional: background, noise, FWHM and ellipticity (3×3 map), background gradient, saturation, trails, sampling}\\
\midrule
FOLDER & folder of stacks (or one image file)\\
--ecrire,\allowbreak{} --write & write QUALITE.csv and QUALITE.txt in each stack\\
--echantillon,\allowbreak{} --sample N & measure N images per stack (spread in time) instead of all\\
--tout,\allowbreak{} --all & measure everything even above 200 images (otherwise: sample of 5 per stack, announced)\\
--processus,\allowbreak{} --processes N & number of measuring processes (default: machine plan)\\
\bottomrule
\end{tabularx}
\par\medskip
\begin{tabularx}{\linewidth}{@{}>{\ttfamily\small\raggedright}p{0.38\linewidth}X@{}}
\toprule
\multicolumn{2}{@{}p{\linewidth}@{}}{\sffamily\bfseries coupole donnees lire (read)\newline\normalfont\small describe the contents of a file (FITS, CSV)}\\
\midrule
FILE & file to read\\
--json & JSON output (for scripts)\\
\bottomrule
\end{tabularx}
\par\medskip
\begin{tabularx}{\linewidth}{@{}>{\ttfamily\small\raggedright}p{0.38\linewidth}X@{}}
\toprule
\multicolumn{2}{@{}p{\linewidth}@{}}{\sffamily\bfseries coupole donnees exporter (export)\newline\normalfont\small export a spectrum or a series as CSV}\\
\midrule
FILE & file to read\\
--csv FILE & export a spectrum or a series as CSV\\
--hdu HDU & index of the dataset in the file (default: the first 1D one)\\
--vitesse,\allowbreak{} --velocity & add the radio velocity (frequency spectrum)\\
--f0 MHz & rest frequency in MHz (default: RESTFRQ, else H I)\\
\bottomrule
\end{tabularx}
\par\medskip
\begin{tabularx}{\linewidth}{@{}>{\ttfamily\small\raggedright}p{0.38\linewidth}X@{}}
\toprule
\multicolumn{2}{@{}p{\linewidth}@{}}{\sffamily\bfseries coupole cosmo\newline\normalfont\small From redshift to distances: comoving, luminosity, angular, lookback time, age, volume, distance modulus, scale (Planck 2018 and other parameter sets)}\\
\midrule
Z & redshift(s) to compute\\
--objet,\allowbreak{} --object NAME & object name: redshift requested from SIMBAD\\
--modele,\allowbreak{} --model \{planck18,\allowbreak{}planck15,\allowbreak{}wmap9,\allowbreak{}simple,\allowbreak{}perso\} & parameter set (default: planck18)\\
--h0 H0 & H₀ in km s⁻¹ Mpc⁻¹ (custom model)\\
--om OM & total Ωm (custom model)\\
--ok OK & Ωk, curvature (planck18 or perso)\\
--shoes & add the SH0ES comparison (planck18)\\
--sans-incertitudes,\allowbreak{} --no-uncertainties & without uncertainties (faster)\\
--csv FILE & write the results to this CSV file\\
--courbes,\allowbreak{} --curves ZMIN ZMAX N & write the curves (z from ZMIN to ZMAX, N points) to the CSV file given by --csv\\
--json & JSON output (for scripts)\\
\bottomrule
\end{tabularx}
\par\medskip
\begin{tabularx}{\linewidth}{@{}>{\ttfamily\small\raggedright}p{0.38\linewidth}X@{}}
\toprule
\multicolumn{2}{@{}p{\linewidth}@{}}{\sffamily\bfseries coupole sites\newline\normalfont\small Observing sites (OpenStreetMap map), time zones, UTC and local time}\\
\midrule
--ajouter,\allowbreak{} --add X & add a site: ID NAME LAT LON ALT TIMEZONE [MPC]\\
--supprimer,\allowbreak{} --delete ID & delete an added site\\
--heure,\allowbreak{} --time ISO\_DATE & show a UTC date (ISO) in site time and evening date\\
--site ID & site identifier (default: ohp)\\
\bottomrule
\end{tabularx}
\par\medskip
\begin{tabularx}{\linewidth}{@{}>{\ttfamily\small\raggedright}p{0.38\linewidth}X@{}}
\toprule
\multicolumn{2}{@{}p{\linewidth}@{}}{\sffamily\bfseries coupole machine\newline\normalfont\small Hardware diagnosis, chosen parallelism, graphics card, ASTAP}\\
\midrule
--json & JSON output (for scripts)\\
\bottomrule
\end{tabularx}
\par\medskip
}

% =====================================================================
\section{Evolution}\label{chap:evolutif}
\paragraph{Sources.} Addresses (TAP, collection, portal, download URL rewriting, resolvers, collection point,
repository, map tiles) are in \texttt{coupole/donnees/sources.json}. Three layers: shipped values; file published in
the repository, fetched at startup when newer and passing the shape check (known keys, types, domains from an
allow-list; no cryptographic signature); values forced by the user, never overwritten. Preferences > Sources: edit,
back to default, test. If the image server moves, \texttt{ohp.telechargement.de}/\texttt{vers} rewrites the URLs.
\paragraph{Classification.} Shipped table \texttt{data/classement.json} (194 checked names); published aliases; user
corrections (\emph{Correct the classification}, \cmd{coupole ohp classify}); name-shape rules; optional online
resolution (JPL SBDB, CDS Sesame), cached; grouping by position on a known fixed object; otherwise « unclassified ».
Every automatic classification is marked « to check ».
\paragraph{Database.} Extra ObsCore columns ignored, missing ones replaced by a default, unknown instruments and
filters stored apart, non-image data never flagged « daytime »; tested on simulated inventories.
\paragraph{Security.} No secret in the repository; remote sources file and collection point over HTTPS only, bounded
size, domain allow-list; update archive downloaded from GitHub only and checked before extraction (absolute paths or
« .. » refused, symbolic links refused, unpacked size and member count bounded), replacement with rollback on
failure; XISF and FITS read with size bounds; folder names built from the database made safe (ASCII, forbidden
characters, Windows reserved names, length); ASTAP launched without a shell; no \texttt{eval} nor \texttt{pickle}
on external data.

% =====================================================================
\section{Incident reports and privacy}
At first launch Coupole asks whether it may send anonymous reports: installation, sign of life (once a day), crash,
native crash, frozen interface, written report. Content: version, system, processor, memory, graphics card, a random
installation identifier, error trace with the home folder replaced by « \textasciitilde », the user name and the
computer name masked and deep paths shortened; a report never exceeds 64 kB. Never a computer name, a user name or an
image. Without consent nothing is sent: reports stay
in \texttt{\_rapports} in the settings folder. The choice can be changed (Help menu, \cmd{coupole reports
--enable|--disable}).

% =====================================================================
\section{Writing a module}
A module is a Python package with a \texttt{manifest.json} (identifier, version, French and English names, icon,
texts, graphical and command-line entry points), discovered in \texttt{coupole/modules/}, through the
\texttt{coupole.modules} entry point, or in the \texttt{modules/} folder of the settings. Texts through \texttt{tr()},
mandatory tooltips, long work off the GUI thread: see CONTRIBUTING.md.

% =====================================================================
\section{Troubleshooting}
\begin{center}\small
\begin{tabularx}{\linewidth}{@{}>{\raggedright}p{0.34\linewidth}X@{}}
\toprule
\textbf{Symptom} & \textbf{What to do}\\
\midrule
Nothing starts (Windows) & Start menu > Coupole > Diagnostic: the error stays on screen.\\
« Graphical interface unavailable » & missing Qt libraries (Linux: \texttt{libxcb-cursor0}, \texttt{libegl1}); the command line still works.\\
Download fails & run again (resume); check the connection; Preferences > Sources > Test; server moved: URL rewriting.\\
« ASTAP not found » & ASTAP assistant; or \texttt{coupole astap --set PATH --catalog FOLDER}.\\
Incomplete catalogue & reinstall the catalogue (1,476 tiles expected).\\
Not enough space & choose another output folder; the space needed includes pending FITS.\\
Slow computer & economy mode, or fewer simultaneous conversions.\\
Unknown target name & Catalogue tab, « To check » box, \emph{Correct the classification}.\\
Quality unavailable & \texttt{pip install sep} (included in the packages).\\
Settings back to defaults & the \texttt{reglages.json} file was unreadable: it is kept alongside with the suffix \texttt{.corrompu-…}.\\
Output folder not writable & clear message at start; choose another folder (Preferences or the Folder field).\\
\texttt{.part} or \texttt{.tmp} file in \texttt{\_traitement} & left by a cut: resuming reuses it (\texttt{.part}) or replaces it; nothing to do.\\
\bottomrule
\end{tabularx}
\end{center}

\section*{Sources}
\addcontentsline{toc}{section}{Sources}
« OHP student observations » bank, Observatoire de Paris / PADC (Paris Astronomical Data Centre), public TAP
service, Etalab 2.0 licence; official guide of the database (PADC-BDD-DU-ECU). XISF 1.0 specification (Pleiades
Astrophoto). ASTAP: \url{https://www.hnsky.org/astap.htm}. Observatory codes: Minor Planet Center. SEP: Barbary
(2016), Bertin \& Arnouts (1996). Adaptive moments: Bernstein \& Jarvis (2002). Map background: © OpenStreetMap
contributors. SIMBAD and Sesame: CDS, Strasbourg (Wenger et al. 2000); Aladin Lite (CDS); NED (NASA/IPAC); JPL
Small-Body Database (NASA/JPL). Cosmology: Planck Collaboration (2020, A\&A 641, A6); Hinshaw et al. (2013, WMAP 9);
Riess et al. (2022, SH0ES); Hogg (1999, astro-ph/9905116); Komatsu et al. (2011) for the massive-neutrino density;
astropy.cosmology; « cosmologie-redshift » calculator (same author).
\end{document}
